\documentclass[10pt,aspectratio=169]{beamer}
\usetheme{gotham}
\input{gotham-personal}

% ── Packages ──────────────────────────────────────────────────────────────────
\usepackage{booktabs}
\usepackage{graphicx}
\usepackage{amsmath, amssymb}
\usepackage[numbers,sort&compress]{natbib}
\usepackage{hyperref}
\usepackage{tikz}
\usetikzlibrary{arrows.meta, positioning}

% ── Metadata ──────────────────────────────────────────────────────────────────
\title[Identity norms \& work]{Identity norms and the decision to work}
\subtitle{The income-elastic breadwinner norm}
\author[Saggese]{Allegra Saggese\textsuperscript{1}}
\institute{\textsuperscript{1}UC Santa Cruz}
\date[May 2026]{May 2026}

% ── Convenience macros ────────────────────────────────────────────────────────
\newcommand{\fig}[2][0.82\textwidth]{%
  \begin{center}
    \includegraphics[width=#1]{#2}
  \end{center}
}
\newcommand{\twofig}[2]{%
  \begin{columns}[c]
    \column{0.5\textwidth}\centering\includegraphics[width=\linewidth]{#1}
    \column{0.5\textwidth}\centering\includegraphics[width=\linewidth]{#2}
  \end{columns}
}

% ── Graph path root (relative to this file) ───────────────────────────────────
\graphicspath{{../data/graphs/}}

% ── Citation style ────────────────────────────────────────────────────────────
\setbeamertemplate{bibliography item}[text]

% ==============================================================================
\begin{document}
% ==============================================================================

\maketitle

% ── Agenda ────────────────────────────────────────────────────────────────────
\begin{frame}[toc]{Agenda}
  \tableofcontents
\end{frame}

% ==============================================================================
\section{Motivation}
% ==============================================================================

\begin{frame}{The long-run picture}
  \begin{columns}[c]
    \column{0.60\textwidth}
      % Replace with FRED screenshot: LFPR Women (green) and Men (purple), 1950-2024
      % Use the same figure as Round 2 slide 3 (FRED export)
      \fbox{\parbox{\linewidth}{\centering\vspace{2cm}
        \textit{[Figure: Male and Female LFPR, US, 1950–2024 \\
        Source: FRED, BLS]}
      \vspace{2cm}}}
    \column{0.40\textwidth}
      \begin{itemize}
        \item Female LFPR rose from 33\% (1950) to $\sim$58\% (2024)
        \item \textbf{Since $\sim$2000: gains have stalled}
        \item Male LFPR declining; female LFPR flat
        \item No country has female LFPR exceeding male
              \citep{SaezGethin2026, RameyFrancis2009}
      \end{itemize}
  \end{columns}
  \vspace{0.5em}
  \begin{block}{The puzzle}
    If institutional barriers have fallen \citep{Goldin1992, Goldin2023},
    why has progress stalled?
  \end{block}
\end{frame}

\begin{frame}{Where women have stopped gaining}
  \begin{columns}[c]
    \column{0.55\textwidth}
      \begin{itemize}
        \item Female gains came primarily on the \textbf{extensive margin}
              (more women entering) \citep{SaezGethin2026}
        \item Today: US women now exceed men in college attainment
              (47\% BA/BS vs.\ 37\% for men; Pew 2024)
        \item Women earn more than men in 22 metro areas (Pew 2022)
        \item And yet: \textbf{no country has seen female LFPR cross male LFPR}
      \end{itemize}
    \column{0.45\textwidth}
      \begin{exampleblock}{US is anomalous}
        Women's participation still growing in most OECD countries.
        US stagnated earlier and is farther behind — not explained
        by religion or institutions alone.
      \end{exampleblock}
      \medskip
      \begin{alertblock}{Implication}
        The barrier today is not a broken ladder or glass ceiling
        preventing entry — it constrains \emph{how much} women work.
      \end{alertblock}
  \end{columns}
\end{frame}

\begin{standinenv}
\begin{frame}{Research question}
  Do persistent gender norms materialize in women's US labor market
  decisions today — and is the norm stronger where households
  can \emph{afford} to enforce it?
\end{frame}
\end{standinenv}

% ==============================================================================
\section{What we know}
% ==============================================================================

\begin{frame}{How women entered the labor force}
  Recall the evidence on the rise in female LFPR (extensive margin),
  adapted from \citet{Goldin1992, Goldin2023}:
  \begin{itemize}
    \item Development of birth control
    \item Removal of legislation prohibiting women from teaching and clerical jobs
    \item Title IX and women's workplace protections
    \item The Equal Pay Act and Civil Rights Act
    \item Rising access to education; mainstreaming of the dual-income household
  \end{itemize}
  \medskip
  \textbf{Today:} institutional barriers have largely fallen.
  \textit{Something else} is now holding the intensive margin down.
\end{frame}

\begin{frame}{Evidence is building that norms play a role}
  \begin{itemize}
    \item[$\Rightarrow$] Norms have been neglected in economics because they are
          incompatible with modeling preferences cleanly \citep{Lundberg2025}
    \item[$\Rightarrow$] Gender norms influence Japanese households to shift female
          labor toward domestic work to minimize household cost
          \citep{SakamotoKohara2025}
    \item[$\Rightarrow$] Gender norms constrain economic growth in developing
          countries, like a poverty trap \citep{Fieldetal2026}
    \item[$\Rightarrow$] Historical evidence from America's frontier shows the
          \textbf{persistence of gender norms} in areas that were once
          homesteader populations \citep{Bazzietal2025}
    \item[$\Rightarrow$] Masculinity norms persist when the sex ratio favors men
          \citep{Baranovetal2023}
  \end{itemize}
\end{frame}

\begin{frame}{The breadwinner norm: Bertrand, Kamenica \& Pan (2015)}
  \citet{BertrandKamenicaPan2015} document that \textbf{US couples
  systematically avoid the wife out-earning the husband}:
  \begin{columns}[c]
    \column{0.48\textwidth}
      \begin{itemize}
        \item Sharp drop in couple density at $z = 0.5$ (wife's share of
              household earnings)
        \item When wife earns more: lower marriage rates, higher divorce,
              lower female LFPR
        \item Effect robust across demographics and time periods
              (CPS/Census, 1970–2013)
        \item \textbf{Reported uniformly across the income distribution}
      \end{itemize}
    \column{0.52\textwidth}
      % BKP Table II screenshot from Round 2 slides
      \fbox{\parbox{\linewidth}{\centering\vspace{1.8cm}
        \textit{[BKP Table II — \\
        Wife in labor force $\sim$ PrWifeEarnsMore]}
      \vspace{1.8cm}}}
  \end{columns}
  \medskip
  \begin{alertblock}{Gap in BKP}
    They do not test whether the threshold effect varies with
    \emph{household income level}. We fill this gap.
  \end{alertblock}
\end{frame}

% ==============================================================================
\section{A new narrative}
% ==============================================================================

\begin{frame}{A five-point historical arc}
  \begin{enumerate}
    \item Women gained exposure to work to fulfill a patriotic duty under WWII
    \item Returned home under booming postwar conditions (rising GDP, low
          inequality, government encouragement of homeownership)
    \item Institutional transformations (Goldin 1992) and rising opportunity
          cost of having children drove female LFPR up from the 1970s onward
    \item 1980s–2000s: rising wage inequality forced \textbf{dual-income necessity}
          at middle incomes — economic constraints overrode norm at Q1–Q3
    \item Today: top incomes have exploded; at the top quintile, the ``price''
          of norm enforcement (foregone female income) is once again affordable
  \end{enumerate}
  \medskip
  \textbf{Key implication:} the norm did not change — \emph{who can act on it did.}
\end{frame}

\begin{frame}{Wealth inequality re-stratified norm enforcement}
  \begin{columns}[c]
    \column{0.50\textwidth}
      \begin{block}{1950s–1960s (low inequality)}
        Breadwinner arrangement broadly affordable.\\
        Norm had roughly uniform bite across middle class.
      \end{block}
      \medskip
      \begin{block}{1980s–2000s (rising inequality)}
        Dual-income becomes economic necessity for Q1–Q3.\\
        Economics override the norm at middle incomes.\\
        \textbf{Norm appears to go dormant.}
      \end{block}
      \medskip
      \begin{alertblock}{2010s–2020s (high inequality)}
        Top quintile income explodes.\\
        Foregone wife's income ($\sim$\$50–60k) trivial at \$150k+ male income.\\
        \textbf{Norm re-emerges — now concentrated at the top.}
      \end{alertblock}
    \column{0.50\textwidth}
      \begin{itemize}
        \item BKP (2015) data run to 2013 — they find a uniform average effect
        \item That average is dominated by Q4–Q5 households
        \item The income-elastic pattern would be invisible if not stratified
        \item \textbf{New prediction:} BKP avoidance behavior should be
              increasingly top-concentrated as inequality rises
        \item Testable in 1980–2023 IPUMS data: run BKP density ratio by
              income decile and year
      \end{itemize}
  \end{columns}
\end{frame}

% ==============================================================================
\section{The income-elastic norm: stylized facts}
% ==============================================================================

\begin{frame}{Stylized fact 1 — the gap is income-elastic}
  \textbf{Female working share by HH income quintile, Democrat vs.\ Republican counties,
  IPUMS 2010–2020}
  \fig{2026-05-07_ipums_hh_gradient_dem_vs_rep_2010_2020.png}
  \begin{itemize}
    \item Q1–Q3: \textbf{zero gap} between Dem and Rep counties
    \item Q4: $+0.9$pp gap; \textbf{Q5: $+2.35$pp gap}
    \item Both groups peak at Q4 then decline at Q5; Rep decline is steeper ($-3.3$pp vs.\ $-1.9$pp)
  \end{itemize}
\end{frame}

\begin{frame}{Stylized fact 2 — male-only earner re-emerges at the top}
  \textbf{Share of male-only earner households by income quintile,
  Democrat vs.\ Republican counties}
  \fig{2026-05-07_ipums_hh_male_only_gradient_dem_vs_rep_2010_2020.png}
  \begin{itemize}
    \item At Q5, Republican counties: male-only earner rises from 16\% $\to$ 20\%
    \item Democrat counties Q5: 16\% $\to$ 18\%
    \item \textbf{The traditional breadwinner household is a high-income conservative
          phenomenon, not a poverty story}
  \end{itemize}
\end{frame}

\begin{frame}{Stylized fact 3 — the norm is entirely on the wife's side}
  \textbf{Weekly hours by income decile — husband (dashed) and wife (solid),
  Democrat vs.\ Republican counties}
  \fig[0.78\textwidth]{2026-05-11_hours_by_income_decile_husband_wife_dem_vs_rep.png}
  \begin{columns}[c]
    \column{0.5\textwidth}
      \begin{itemize}
        \item Wife hours peak D8–D9, then \textbf{drop at D10} for both groups
        \item Rep wife D10: 29.2 hrs/wk; Dem wife D10: 31.4 hrs/wk
              ($-2.2$ hr gap at the top)
      \end{itemize}
    \column{0.5\textwidth}
      \begin{itemize}
        \item \textbf{Husband hours: monotonically rising, nearly identical across
              political groups} (Rep 45.9, Dem 46.2 at D10)
        \item The pullback is gendered and concentrated at the top
      \end{itemize}
  \end{columns}
\end{frame}

\begin{frame}{Stylized fact 4 — gap requires a strong political signal}
  \textbf{Q5 county female LFPR by political lean}
  \fig[0.70\textwidth]{2026-05-07_departure_q5_female_lfpr_by_political_lean.png}
  \begin{itemize}
    \item Among nationally Q5 counties:
          Strongly Dem 76.8\%, Lean Dem 75.7\%, Lean Rep 76.6\%,
          \textbf{Strongly Rep 74.0\%}
    \item Gap is \textbf{non-monotonic}: Lean Rep has \emph{higher} LFPR than Lean Dem
    \item Norm only operates with $>10$pp vote margin — not a smooth cultural gradient
    \item Consistent with an identity norm that requires social reinforcement
          to operate \citep{AkerlofKranton2000}
  \end{itemize}
\end{frame}

\begin{frame}{Stylized fact 5 — the gap is persistent, not converging}
  \textbf{Q5 county female LFPR trend, Democrat vs.\ Republican, 2012–2020}
  \fig[0.72\textwidth]{2026-05-07_departure_q5_female_lfpr_trend_dem_vs_rep.png}
  \begin{itemize}
    \item Dem counties Q5: stable $\sim$76–77\%; Rep counties Q5: stable $\sim$74–74.5\%
    \item Gap is \textbf{flat}: not a legacy of earlier sorting that is converging away
    \item Consistent with a persistent norm, not a transitional composition effect
  \end{itemize}
\end{frame}

% ==============================================================================
\section{Evidence for the mechanism}
% ==============================================================================

\begin{frame}{Her threat point weakens exactly where the norm operates}
  \textbf{Spousal income (no transfers) vs.\ log(HH income), by anchor year}
  \fig[0.82\textwidth]{2026-05-07_ipums_spouse_wage_income_by_hh_income_anchor_years.png}
  \begin{itemize}
    \item At \$200k HH income (2010): husband earns $\sim$\$150k, wife $\sim$\$50–60k
    \item \textbf{Male income: near-linear with HH income}
    \item \textbf{Female income: plateaus} — her contribution flattens as HH income rises
    \item Pattern consistent and sharpening 1980–2023
    \item The husband's outside option grows; hers does not — bargaining power shifts
  \end{itemize}
\end{frame}

\begin{frame}{Wife's earned share by income decile}
  \textbf{Wife's share of couple earned income, Democrat vs.\ Republican counties}
  \fig[0.72\textwidth]{2026-05-11_share_vs_hh_income_decile_dem_vs_rep.png}
  \begin{itemize}
    \item Wife's share falls monotonically D1 $\to$ D10 for both groups
          ($\sim$42\% $\to$ $\sim$33\%)
    \item Gap at D10: Dem 34.3\%, Rep 32.2\% $\Rightarrow$ \textbf{2.08pp gap}
    \item Gap is small and parallel D1–D9; \textbf{widens sharply at D10}
    \item Consistent with norm suppressing her income share precisely where husband's
          outside option is strongest
  \end{itemize}
\end{frame}

\begin{frame}{Avoidance behavior at the breadwinner threshold}
  \textbf{Density of wife's earned income share ($z_{\text{earned}}$),
  Democrat vs.\ Republican counties — donut $\pm 2$pp at $z = 0.5$}
  \fig[0.75\textwidth]{2026-05-11_rdd_donut_density_earned_income_share_dem_vs_rep.png}
  \begin{columns}[c]
    \column{0.5\textwidth}
      \begin{itemize}
        \item Clear bunching just below $z = 0.5$ in both groups
        \item Below/above $z=0.5$ ratio: \\
              Dem: $14.5\% / 10.2\% = \mathbf{1.42}$ \\
              Rep: $14.4\% / 9.6\% = \mathbf{1.50}$
      \end{itemize}
    \column{0.5\textwidth}
      \begin{itemize}
        \item \textbf{Rep counties show $\sim$5\% stronger avoidance}
        \item Replicates BKP (2015) nationally; extends to Dem/Rep heterogeneity
        \item Donut design removes mechanical mass point at exactly $z=0.50$
              (3.96\% of weighted obs)
      \end{itemize}
  \end{columns}
\end{frame}

\begin{frame}{Kink RDD: wife's hours collapse at the breadwinner threshold}
  \textbf{Kink RDD on wife's weekly hours — running variable $z_{\text{earned}}$,
  donut $\pm 2$pp, bandwidth $\pm 20$pp}
  \fig[0.75\textwidth]{2026-05-11_rdd_donut_kink_wife_hours_earned_income.png}
  \begin{columns}[c]
    \column{0.55\textwidth}
      \begin{itemize}
        \item Below threshold: wife scales up hours $+27.5$ hrs/pp as share rises
        \item \textbf{At threshold: slope nearly collapses} — net slope $\sim 2.3$ hrs/pp
        \item Kink = $-25.2$ hrs/pp ($p < 0.001$)
        \item Jump = $-0.386$ hrs ($p < 0.001$)
        \item Wife employment: Jump $= -0.007$ ($p < 0.001$)
      \end{itemize}
    \column{0.45\textwidth}
      \begin{itemize}
        \item \textbf{Husband hours: NOT kinked} (Jump $= +0.12$, symmetric kink)
        \item Robustness: kink stable at $\sim$25 hrs/pp across donut widths
              $\pm 1$pp, $\pm 2$pp, $\pm 3$pp
        \item Identification: she is already working and has positive earned
              income — pure wealth effects do not explain the avoidance at the
              threshold
      \end{itemize}
  \end{columns}
\end{frame}

% ==============================================================================
\section{Model}
% ==============================================================================

\begin{frame}{Departure from existing models}
  \begin{columns}[c]
    \column{0.50\textwidth}
      \begin{block}{Existing norm models}
        \citet{AkerlofKranton2000}, \citet{FernandezFogliGuner2005},
        \citet{FogliVeldkamp2011}, \citet{BisinVerdier2000}:
        \begin{itemize}
          \item Norm strength $\tau$ is a fixed cultural parameter
          \item Transmitted across generations; varies by county or group
          \item Does not interact with household income
        \end{itemize}
        \textbf{Prediction:} uniform gap across income distribution
        within conservative counties.
      \end{block}
    \column{0.50\textwidth}
      \begin{alertblock}{Our departure}
        Norm penalty is \textbf{income-elastic}:
        $\theta(\tau, y)$ increasing in both norm intensity $\tau$
        and household income $y$.
        \smallskip

        Interpretation: a richer conservative household can
        \emph{afford} to enforce the norm — the foregone income
        of the wife is small relative to total household resources.
        \smallskip

        \textbf{Prediction:} zero gap at low income; positive and
        growing gap at high income — exactly what we observe.
      \end{alertblock}
  \end{columns}
\end{frame}

\begin{frame}{Model: joint utility with income-elastic norm penalty}
  \textbf{Setup:} Two agents $(H, W)$; household jointly maximizes utility.
  \begin{align*}
    U &= u(c) - v(l_H) - v(l_W) - \theta(\tau, y) \cdot \mathbf{1}\!\left[\frac{y_W}{y_H} \geq 1\right]
  \end{align*}
  where $c$ = consumption, $l_i$ = labor hours, $y_i$ = individual earned income,
  $\tau$ = cultural norm intensity (county political lean), $y = y_H + y_W$.

  \bigskip
  \textbf{Norm penalty:}
  \[
    \theta(\tau, y) = \underbrace{\tau}_{\text{norm intensity}} \times \underbrace{g(y)}_{\text{increasing in }y}
  \]
  At low income: $g(y) \approx 0$ $\Rightarrow$ economics dominate $\Rightarrow$ zero gap. \\
  At high income: $g(y)$ large $\Rightarrow$ norm operates $\Rightarrow$ positive gap.

  \bigskip
  \textbf{Enforcement:} Both spouses internalize $\theta$
  \citep{AkerlofKranton2000} — norm is not an external constraint but a
  shared identity cost, reinforced by in-group culture.
\end{frame}

\begin{frame}{Model predictions vs.\ data}
  \begin{columns}[c]
    \column{0.52\textwidth}
      \begin{block}{Model predictions}
        \begin{enumerate}
          \item Zero gap at Q1–Q3: $\theta(\tau, y_{\text{low}}) \approx 0$
          \item Positive gap at Q4–Q5: $\theta(\tau, y_{\text{high}}) > 0$,
                increasing in $\tau$
          \item Wife's hours kink at $z = 0.5$, sharper in high-$\tau$ counties
          \item Husband's hours: unaffected by $\tau$ or income (his outside
                option grows, but the identity cost falls on the wife)
          \item Gap strengthens over time as top-quintile income rises
                (via $g(y)$)
        \end{enumerate}
      \end{block}
    \column{0.48\textwidth}
      \begin{exampleblock}{Calibration targets}
        \begin{tabular}{ll}
          \toprule
          Moment & Value \\
          \midrule
          Work share gap Q1 & 0pp \\
          Work share gap Q5 & 2.35pp \\
          Hours gap D10 (Dem) & 14.8 hrs \\
          Hours gap D10 (Rep) & 16.7 hrs \\
          BKP ratio (Dem)  & 1.42 \\
          BKP ratio (Rep)  & 1.50 \\
          Kink hrs/pp      & $-25.2$ \\
          \bottomrule
        \end{tabular}
      \end{exampleblock}
  \end{columns}
\end{frame}

% ==============================================================================
\section{Next steps}
% ==============================================================================

\begin{frame}{Identification strategy}
  \textbf{Causal identification: school board elections (Ballotopedia)}
  \begin{itemize}
    \item Descriptive evidence is consistent with norm suppression, but
          cannot rule out time-invariant sorting or industrial composition
    \item \textbf{Planned event study:} close elections to conservative school
          boards shift county-level gender norm environment
    \item Outcome: female LFPR, wife's hours, wife's earned share — at the
          household level, stratified by income decile
    \item \textbf{Expected pattern under income-elastic norm hypothesis:}
          a conservative school board win should have zero effect at Q1–Q3
          and a negative effect on wife's hours/LFPR at Q4–Q5
  \end{itemize}
  \medskip
  \textbf{Additional robustness:}
  \begin{itemize}
    \item RDD political interaction: is the kink in wife's hours sharper
          in strongly Republican counties? (designed, not yet run)
    \item Better norm proxy: CCES county-level gender attitudes (direct survey
          items on breadwinner expectations), ARDA evangelical share
    \item Occupation/education controls to rule out labor market composition
  \end{itemize}
\end{frame}

\begin{frame}{Takeaways}
  \begin{enumerate}
    \item \textbf{New stylized fact:} the conservative-liberal LFPR differential
          is zero below the top income quintile and concentrates sharply at Q5.
          This income-elastic pattern has not been previously documented.
    \smallskip
    \item \textbf{The norm is on the wife, not the husband:} at D10, wives in
          Republican counties work 2.2 fewer hours per week than in Democratic
          counties; husbands are virtually identical.
    \smallskip
    \item \textbf{Breadwinner threshold:} replicating and extending BKP (2015),
          a kink of $-25.2$ hrs/pp at $z=0.5$ (wife's earned share), stronger
          avoidance in Republican counties (ratio 1.50 vs.\ 1.42).
    \smallskip
    \item \textbf{Model departure:} existing models treat norm strength as fixed.
          We propose an income-elastic norm penalty $\theta(\tau, y)$, consistent
          with rising wealth inequality re-stratifying who can enforce the norm.
  \end{enumerate}
\end{frame}

% ── Questions ─────────────────────────────────────────────────────────────────
\begin{frame}[standout]
  Questions?
\end{frame}

% ==============================================================================
% References
% ==============================================================================
\begin{frame}[allowframebreaks]{References}
  \footnotesize
  \begin{thebibliography}{99}
    \bibitem{AkerlofKranton2000}
      Akerlof, G.\ A., \& Kranton, R.\ E.\ (2000).
      Economics and identity.
      \textit{Quarterly Journal of Economics}, 115(3), 715--753.
    \bibitem{Baranovetal2023}
      Baranov, V., et al.\ (2023).
      Masculinity norms and economic behavior.
      \textit{Working paper}.
    \bibitem{Bazzietal2025}
      Bazzi, S., et al.\ (2025).
      The frontier of gender norms.
      \textit{Working paper}.
    \bibitem{BertrandKamenicaPan2015}
      Bertrand, M., Kamenica, E., \& Pan, J.\ (2015).
      Gender identity and relative income within households.
      \textit{Quarterly Journal of Economics}, 130(2), 571--614.
    \bibitem{BisinVerdier2000}
      Bisin, A., \& Verdier, T.\ (2000).
      ``Beyond the melting pot'': Cultural transmission, marriage, and
      the evolution of ethnic and religious traits.
      \textit{Quarterly Journal of Economics}, 115(3), 955--988.
    \bibitem{FernandezFogliGuner2005}
      Fern\'{a}ndez, R., Fogli, A., \& Olivetti, C.\ (2004).
      Mothers and sons: Preference formation and female labor force dynamics.
      \textit{Quarterly Journal of Economics}, 119(4), 1249--1299.
    \bibitem{Fieldetal2026}
      Field, E., et al.\ (2026).
      Gender norms and poverty traps.
      \textit{Working paper}.
    \bibitem{FogliVeldkamp2011}
      Fogli, A., \& Veldkamp, L.\ (2011).
      Nature or nurture? Learning and the geography of female labor force
      participation.
      \textit{Econometrica}, 79(4), 1103--1138.
    \bibitem{Goldin1992}
      Goldin, C.\ (1992).
      Understanding the gender gap: An economic history of American women.
      \textit{Oxford University Press}.
    \bibitem{Goldin2023}
      Goldin, C.\ (2023).
      Career and family: Women's century-long journey toward equity.
      \textit{Princeton University Press}.
    \bibitem{Lundberg2025}
      Lundberg, S.\ (2025).
      Norms and economics.
      \textit{Working paper}.
    \bibitem{RameyFrancis2009}
      Ramey, V.\ A., \& Francis, N.\ (2009).
      A century of work and leisure.
      \textit{American Economic Journal: Macroeconomics}, 1(2), 189--224.
    \bibitem{SakamotoKohara2025}
      Sakamoto, A., \& Kohara, M.\ (2025).
      Gender norms and female labor allocation.
      \textit{Working paper}.
    \bibitem{SaezGethin2026}
      Saez, E., \& Gethin, A.\ (2026).
      Labor force participation trends.
      \textit{Working paper}.
  \end{thebibliography}
\end{frame}

\end{document}
